\section{Finite-field square constraints}
\label{sec:finite-field-graphs}

This section controls systems of square constraints on independent field
variables.  A single constraint is close to a constant kernel in operator
norm.  Replacing constraints one at a time then counts any fixed directed
graph without loops or opposite edges.  Applying this count to a four-cycle
of two lists gives the operator estimate needed for the tuple measures.

Throughout, finite sets carry uniform probability measure, and
\(e(t)=\exp(2\pi i t)\).  For an odd prime \(p\), write
\[
 Q_p^*=\{u^2:u\in\mathbb F_p^\times\},
 \qquad
 A_p(t)=1_{Q_p^*}(t),
 \qquad
 \delta_p=\frac{\sqrt p+1}{2p}\le p^{-1/2}.
\]
For finite sets \(X,Y\), the operator associated with a kernel
\(K:X\times Y\to\mathbb R\) is
\[
 T_K f(x)=\mathbb E_{y\in Y}K(x,y)f(y).
\]
We use \(\|K\|_{\mathrm{op}}\) to mean the norm of \(T_K\) from
\(L^2(Y)\) to \(L^2(X)\).  A constant kernel \(c\) therefore denotes
the operator \(f\mapsto c\,\mathbb E f\).

\begin{lemma}[Counting directed square constraints]
\label{lem:graph-count}
Let \(p\) be an odd prime and let \(G=(W,E)\) be a finite directed graph
with no loops and at most one edge on each unordered pair of distinct
vertices.  Put \(m=|E|\), and let \((z_w)_{w\in W}\) be independent
uniform elements of \(\mathbb F_p\).  Then
\[
 \left|
 \mathbb E\prod_{(u,v)\in E}A_p(z_v-z_u)-2^{-m}
 \right|
 \le m\delta_p.
\]
The assertion includes the empty graph, whose product is \(1\).
\end{lemma}

\begin{proof}
We first show that the kernel
\(b_p(x,y)=A_p(y-x)-1/2\) has operator norm at most \(\delta_p\).
The characters \(x\mapsto e(kx/p)\), \(k\in\mathbb F_p\), are an
orthonormal eigenbasis.  The constant eigenvalue is
\[
 \frac{|Q_p^*|}{p}-\frac12=-\frac1{2p}.
\]
For \(k\ne0\), put \(G_p(k)=\sum_{u\in\mathbb F_p}e(ku^2/p)\).
The corresponding eigenvalue is
\[
 \frac1p\sum_{t\in Q_p^*}e(kt/p)
 =\frac{G_p(k)-1}{2p}.
\]
Here the subtraction of \(1\) removes the zero square, and division
by \(2\) removes the multiplicity of every nonzero square.  Since
\((u,v)\mapsto(u-v,u+v)\) is a bijection on \(\mathbb F_p^2\),
\[
 |G_p(k)|^2
 =\sum_{u,v\in\mathbb F_p}e\bigl(k(u^2-v^2)/p\bigr)
 =\sum_{a,b\in\mathbb F_p}e(kab/p)
 =p.
\]
Thus all eigenvalues have absolute value at most \(\delta_p\).
In particular,
\begin{equation}
 \left|
 \mathbb E_{x,y} f(x)g(y)\bigl(A_p(y-x)-\tfrac12\bigr)
 \right|
 \le\delta_p\|f\|_2\|g\|_2.
 \label{eq:single-square-mixing}
\end{equation}
The same bound holds with \(x-y\) in place of \(y-x\), since the
transposed operator has the same norm.

Now replace the \(m\) edge indicators one at a time by \(1/2\).
In a telescoping error for an edge \((u,v)\), fix every variable
except \(z_u,z_v\).  No other edge has both these endpoints.
Consequently all the other factors separate as
\(C f(z_u)g(z_v)\), where \(|C|,|f|,|g|\le1\).
Equation~\eqref{eq:single-square-mixing} bounds the conditional error
by \(\delta_p\).  Averaging the fixed variables and summing over
the edges proves the assertion.
\end{proof}

We next divide the variables into two independent lists.  Only
constraints between the lists enter the following kernel; constraints
within either list must be treated separately.

\begin{lemma}[Mixing between two lists]
\label{lem:list-mixing}
Let \(I,J\) be finite disjoint sets, and let \(G\) be a directed graph
on \(I\sqcup J\) whose \(m\) edges all join \(I\) to \(J\), with at
most one edge on each unordered pair.  The directions may vary from
edge to edge.  For
\(x\in\mathbb F_p^I\), \(y\in\mathbb F_p^J\), write \(z_i=x_i\) for
\(i\in I\), \(z_j=y_j\) for \(j\in J\), and set
\[
 K(x,y)=\prod_{(u,v)\in E(G)}A_p(z_v-z_u).
\]
Then
\begin{equation}
 \|K-2^{-m}\|_{\mathrm{op}}
 \le (32m)^{1/4}p^{-1/8}.
 \label{eq:list-mixing}
\end{equation}
For \(m=0\), both sides are zero.  In particular, the implied
constant for any fixed pair of lists is independent of \(p\).
\end{lemma}

\begin{proof}
Put \(c=2^{-m}\) and \(B=K-c\).  With respect to the normalized
point-mass bases, \(T_B\) has matrix entries
\(B(x,y)/\sqrt{|\mathbb F_p^I|\,|\mathbb F_p^J|}\).
If \(s_r\) are its singular values, then
\begin{align}
 \|B\|_{\mathrm{op}}^4
 &\le \sum_r s_r^4
 =\operatorname{tr}\bigl((T_BT_B^*)^2\bigr) \notag\\
 &=\mathbb E_{x^0,x^1,y^0,y^1}
 B(x^0,y^0)B(x^0,y^1)B(x^1,y^0)B(x^1,y^1),
 \label{eq:list-four-cycle}
\end{align}
where the four lists are independent and uniform.  The kernel is
real, so no conjugates appear in the last expression.

Let \(\Omega=\{0,1\}^2\).  Expanding the last line gives
\[
 \sum_{F\subseteq\Omega}(-c)^{4-|F|}
 \mathbb E\prod_{(a,b)\in F}K(x^a,y^b).
\]
For each \(F\), its product is the constraint indicator of a graph
on
\[
 (I\times\{0,1\})\sqcup(J\times\{0,1\}).
\]
Every base edge between \(i\in I\) and \(j\in J\) is copied, in its
original direction, between \((i,a)\) and \((j,b)\) for every
\((a,b)\in F\).  The two endpoints determine the base unordered
pair and the corner \((a,b)\).  Hence this graph has no loops or
repeated unordered edges, and has exactly \(m|F|\) edges.
Lemma~\ref{lem:graph-count} therefore gives
\[
 \mathbb E\prod_{(a,b)\in F}K(x^a,y^b)
 =c^{|F|}+E_F,
 \qquad |E_F|\le m|F|\delta_p.
\]
The main terms sum to \((c-c)^4=0\).  The remaining terms have
absolute value at most
\[
 m\delta_p\sum_{t=0}^4\binom4t t c^{4-t}
 =4m(1+c)^3\delta_p
 \le32m p^{-1/2}.
\]
Combining this with \eqref{eq:list-four-cycle} and taking fourth
roots proves \eqref{eq:list-mixing}.
\end{proof}
